\section{Open questions and further directions}
\label{sec:open}

The arithmetic results above resolve Nishibuchi's conjecture throughout its
stated parameter range. The local analyticity and differentiability theorem
is also complete under its stated invertibility hypothesis. What remains is
a set of extensions and sharper classifications, together with the broader
questions raised by the original native seed-calculus program.

\subsection{Singular persistence beyond an invertible solved block}
\label{subsec:singular-open}

The weighted lifting theorem allows the solved Jacobian to be noninvertible
modulo $p$ provided that it is invertible over $\Qp$ with a controlled inverse.
This includes a substantial class of lifts often called singular modulo $p$.
The next problem is to understand Jacobians that remain rank-deficient over
$\Qp$, where the inverse construction is no longer available.

For a fixed polynomial system $F\in\Z[X_1,\ldots,X_d]^e$, a precision $L\geq1$,
and a root $\bar x$ modulo $p^L$, let
\[
\mathcal L_{F,L}(\bar x)
=\{x'\in(\Z/p^{L+1}\Z)^d:
F(x')=0,\ x'\equiv\bar x\pmod{p^L}\}.
\]
The finite obstruction theorem decides whether this set is empty from the
residual digit and reduced Jacobian. When it is nonempty, the theorem also
determines its cardinality and affine kernel structure. The conclusion
concerns one level: vanishing of the obstruction permits the next step,
but does not guarantee that subsequent steps are possible.

For example, $X^2-5$ has the root $1$ modulo $2$, and both $1$ and $3$ lift it
to roots modulo $4$. Every odd square is $1$ modulo $8$, so neither branch
continues to a root modulo $8$. Thus even a successful first correction can
lead to branches that stop at the following level.

We seek structural hypotheses and higher obstruction data that control an
entire singular branch, including its behavior under coefficient and parameter
deformation. One route is to propagate suitable conditions from level to
level. Another is to prove existence at all levels in a form that yields
compatible branches. The finite obstruction calculus supplies the tests for
individual levels, but a persistence theorem needs a further argument linking
them. Any general formulation must allow multiple branches, later
obstructions, and changes in the rank of the Jacobian.

\subsection{Native 13.2: presentation invariance}
\label{subsec:native132}

The original conjecture asks for a canonical finite obstruction packet whose
vanishing is necessary and sufficient for a one-level lift in a finitely
presented arithmetic seed system. For a fixed finite polynomial presentation,
the paper now supplies the precise obstruction
\[
o_{F,\bar x}=[r]\in
\coker\!\left(J_{F,\bar x}:\F_p^d\longrightarrow\F_p^e\right),
\]
where $r$ is the residual digit. We have proved independence from the chosen
integer representative. Invariance under a change of equations or presentation
requires more.

The raw cokernel already reveals the difficulty.
The presentations $F=(X)$ and $G=(X,0)$ describe the same zero locus over every
coefficient ring. At zero their Jacobians are, respectively,
\[
\F_p\xrightarrow{\,1\,}\F_p,
\qquad
\F_p\xrightarrow{\,(1,0)\,}\F_p^2.
\]
Their raw cokernels are $0$ and $\F_p$. Both distinguished obstruction classes
vanish, but the ambient obstruction groups are not isomorphic. The example
rules out the raw cokernel as a presentation-independent object while leaving
open the possibility of an intrinsic obstruction construction.

A precise form of Native 13.2 therefore needs four ingredients:
\begin{enumerate}
\item a category of seed objects, including how the operators, relations,
projectors, weights, and topology enter an object;
\item specified changes of presentation and a proof that they preserve the
actual lift problem;
\item an obstruction object that respects dependencies among the relations,
together with comparison maps or a universal property;
\item a precise meaning of ``finite'' and ``canonical'', and a proof that
vanishing is equivalent to the intended one-step continuation.
\end{enumerate}
For commutative algebra presentations, the relation module offers a natural
way to account for dependencies among listed equations. To extend this
approach to the full seed vocabulary, we would need an encoding that preserves
the lift problem, together with a proof of that preservation. This is one of
the links still needed between the concrete obstruction theorem and Native 13.2.

\subsection{Native 13.3: finite compatibility, completion, and descent}
\label{subsec:native133}

The original proposed form is
\[
X\simeq X_2^{\mathrm{wt}}\times_C\prod_{p>2}'X_p.
\]
It leaves $X$, the distinguished subsets defining the restricted product,
the maps to $C$, the ambient category, and the universal meaning of
``smallest'' unspecified. In particular, one must decide whether $X$ is a
global object or a completed object.

The linear compatibility theorem provides one concrete part of this program.
If two local observations
$J:U\to V$ and $K:U\to W$ are linear maps from the \emph{same} correction
module, define
\[
C_{J,K}=U/(\ker J+\ker K).
\]
The attainable observations $\im J$ and $\im K$ have canonical labels in this
quotient. Equality of labels is equivalent to simultaneous realization by a
single correction. The joint image is exactly the fiber product over
$C_{J,K}$, with the proved linear universal property.
The declarations are
\begingroup\small\raggedright
\lean{Komorebi.SharedCorrection.labels_eq_iff_common},
\lean{Komorebi.SharedCorrection.joint_range_eq_fiber}, and
\lean{Komorebi.SharedCorrection.common_universal}.
\par\endgroup
To apply this result arithmetically, one must construct the common correction
module and the observation maps, then show that the resulting compatibility
condition agrees with the intended arithmetic one.

\paragraph{Why the global/completed distinction is necessary.}
Consider constant sequences with integer initial value. Their global parameter
space is $\Z$, while their local parameter spaces are $\Z_p$. Under the
canonical local maps, the image of the global objects is the diagonal
\[
D=\{(n)_p:n\in\Z\}\subseteq\prod_p\Z_p.
\]
This diagonal is dense: a basic open set specifies finitely many congruences,
and the Chinese remainder theorem provides an integer satisfying them. It is
proper: the tuple with $2$-adic coordinate zero and all odd-adic coordinates
one is not the image of an integer. Indeed an integer zero in $\Z_2$ must be
zero, contradicting its $3$-adic coordinate.

Consequently this canonical global diagonal cannot be an equality locus
\[
\{(a,b):f(a)=g(b)\}
\subseteq \Z_2\times\prod_{p>2}\Z_p
\]
for continuous maps to a Hausdorff space $C$. Such an equality locus is closed,
as the inverse image of the closed diagonal in $C\times C$, whereas $D$ is
proper and dense. This topological argument, retained from the native progress
note, concerns the stated topological interpretation. It has not been
formalized in Lean for this manuscript.

Finite CRT reconstruction leads to the familiar completed object
$\widehat\Z=\prod_p\Z_p$, but does not identify which local tuples come from a
single global integer. A complete form of Native 13.3 must therefore specify
finite-level compatibility, transition maps and completion, and any additional
global descent condition. If recurrence, spectral, or Frobenius information
is to supply that condition, it must be expressed through defined maps or
constraints on the local data, with a precise account of how they enter $C$.

The accompanying \nolinkurl{native/Native_Conjectures_Progress.md} records
further results in restricted settings: an intrinsic finite
$\operatorname{Ext}^1$ obstruction for finitely presented commutative integer
algebras, finite CRT reconstruction for polynomial solution sets, and a
characterization of rectangular finite relations by equality of common
labels. The obstruction result is a stated specialization of a classical
external theorem; the CRT and rectangularity results have prose proofs in
the note. These are preserved mathematical progress, with the restrictions
and source records given there, rather than Lean completions of the general
Native 13.2 and 13.3 statements.

\subsection{Arithmetic refinements}
\label{subsec:arithmetic-open}

Write $U_{N,j}(n)=j!C_{N,j}(n)$ and $U_j=U_{4,j}$, and use
$v_p(0)=\infty$ when discussing valuation bounds. The following classification
questions remain open in this program.

\begin{enumerate}
\item \textbf{Higher dyadic differences.}
Determine the exact value of
\[
v_2\!\left(\Delta_u^kU_j(n)\right),\qquad k\geq2,
\]
with explicit treatment of zero differences. The all-order lower bounds and
the exact first-displacement theorem do not determine this value: subtraction
of two first differences can introduce further cancellation. The first
concrete target is order $2$ in one fixed address.

\item \textbf{Least periods beyond $N=4$.}
For further packet families, determine the least positive $P$ for which
$U_{N,j}(n+P)\equiv U_{N,j}(n)\pmod{p^r}$ for every $n$ in the claimed range.
The general sufficient periods are proved. Their minimality does not follow
from those sufficient bounds, and a period in packet index is a degree shift
of $NP$.

\item \textbf{Exact affine Mahler-kernel valuations.}
For the kernel of Section~\ref{subsec:polyphase-exposition}, determine exact
values of
\[
v_p\bigl(\Theta_{M,s}(r,k)\bigr),\qquad
\Theta_{M,s}(r,k)=[z^r](1+z)^s((1+z)^M-1)^k,
\]
beyond the explicit special families. The minimum valuation of the summands
gives a lower bound. Equality requires a separate argument excluding
cancellation.

\item \textbf{Exact classical dyadic tail values.}
The uniform truncation exponent $4M+3$ and its optimality are proved.
Individual term weights are also known. The valuation of the whole omitted
tail at an arbitrary even row may be higher, so an exact formula for all rows
requires more than taking the minimum of those term weights. Empty tails
must remain exactly zero.
\end{enumerate}

These refinements build on the proved estimates and exact cases, including
the sharp uniform classical dyadic bound and the resolution of Nishibuchi's
conjecture.

\subsection{Historical material and interface work}
\label{subsec:historical-scope}

The original Paper IX, with corrected proof-status labels, survives as a
28-page PDF. Its recovered formalization supports the framework, Euler
coordinates, Mahler differences,
unit coordinates, recurrence packets, and CRT results retained in this
manuscript. The exact original TeX bytes, original figure-generation files, and complete
historical experiment package were not recovered with that PDF.
The reconstructed Annals text is reproduced in Appendix~\ref{app:recovered},
with its separate provenance and local qualifications. Their older
figure and test counts remain historical records; those computations have
not been reproduced for this consolidation.

A reconstructed textual copy of the Annals manuscript was also recovered.
It includes finite-field mode and subgroup inversion, Fourier invariance,
Carmichael and matrix transfer, Amice pullbacks, depleted analytic Euler
branches, and extensive polyphase formulas. The sharp classical dyadic kernel
and truncation claims from that text have since been connected to checked
Lean proofs. The finite polyphase formulas used for explanation here have
short explicit algebraic derivations. The remaining claims in that larger
historical text have not all been formalized.
The source archive retains the text with its provenance, and the statement
map distinguishes these historical claims from the certified claims in the
present manuscript.

Remaining recovery work includes the precise older odd-prime edge identities,
their stated ranges, and the binary digit-sum form of the checked weights.
For each item, the work is to recover the formula with its stated range,
check the proof, and identify any corresponding Lean declaration. Where the
historical source is still missing, the record remains incomplete.

One further task is to connect the older chart constructions to the new
regularity theorem. That theorem applies to polynomials with arbitrary
coefficients in the ambient field $\Qp$ and identifies the analytic
germ of every continuous ambient-field root selection under its hypotheses.
Several older rectangular charts were defined using integer coefficients,
integer centers, or subtype-valued maps. Comparing each such historical
interface explicitly with the new ambient formulation remains to be done.
These comparisons would connect the existing constructions to the proved
analytic implicit-function theorem. General singular persistence remains the
separate problem described in Section~\ref{subsec:singular-open}.
